\documentclass[11pt]{amsart}

\usepackage[T1]{fontenc}
\usepackage{lmodern}
\usepackage{microtype}
\usepackage{amsmath,amssymb}
\usepackage[hidelinks]{hyperref}

\hypersetup{
  pdftitle={Minimum Memoization Need Not Lie in the Glushkov MFN Set}
}

\newtheorem{theorem}{Theorem}

\newcommand{\Gl}{\operatorname{Gl}}
\newcommand{\MFN}{\operatorname{MFN}}

\title[Memoization outside the Glushkov MFN set]
{Minimum Memoization Need Not Lie in the Glushkov MFN Set}

\date{}

\subjclass[2020]{68Q45}
\keywords{regular expressions, Glushkov automata, memoization, ambiguity,
feedback vertex set}

\begin{document}

\begin{abstract}
Berglund, van der Merwe, and le Roux conjectured that a minimum
memoization set for a Thompson or Glushkov automaton can be chosen
inside its Minimum Feedback Node set, equivalently that some ordering
makes $\mathrm{IAR}^{+}$ return a minimum set.  We disprove the
Glushkov assertions.  For every $k\geq2$, the Glushkov automaton of
\[
 \left(\left(a\mid
 \underbrace{a^{+}\mid\cdots\mid a^{+}}_{k\ \mathrm{occurrences}}
 \right)(b\mid c)\right)^{*}
\]
has MFN size $k+1$, but its unique minimum memoization set has size $2$
and is disjoint from MFN.  The full MFN set is its only IDA-eliminating
subset, so $\mathrm{IAR}^{+}$ returns all MFN states for every ordering.
Both gaps are unbounded.  The Thompson assertions are not addressed.
\end{abstract}

\maketitle

Berglund, van der Merwe, and le Roux introduced the Minimum Feedback
Node scheme and conjectured that some minimum memoization set is
contained in MFN, equivalently that some ordering of MFN makes
$\mathrm{IAR}^{+}$ return a minimum set~\cite[Conjectures~1 and~2]{Berglund2026}.
We refute the Glushkov clauses of both conjectures.

For an automaton $A$, let $\mu(A)$ be the minimum number of memoized
states needed to eliminate infinite degree of ambiguity (IDA), and set
\[
 \mu_{\mathrm{MFN}}(A)=
 \min\bigl\{|M|:M\subseteq\MFN(A)\text{ and }M\text{ eliminates IDA}\bigr\}.
\]

\begin{theorem}\label{thm:main}
For every $k\geq2$, let
\[
 R_k=\left(\left(a\mid
 \underbrace{a^{+}\mid\cdots\mid a^{+}}_{k\ \mathrm{occurrences}}
 \right)(b\mid c)\right)^{*},
 \qquad G_k=\Gl(R_k).
\]
Then $\MFN(G_k)$ is the unique minimum directed feedback vertex set of
$G_k$, and
\[
 \mu(G_k)=2,
 \qquad
 \mu_{\mathrm{MFN}}(G_k)=|\MFN(G_k)|=k+1.
\]
The unique minimum memoization set is disjoint from $\MFN(G_k)$, and
$\mathrm{IAR}^{+}$ returns all $k+1$ MFN states for every ordering.
\end{theorem}

\begin{proof}
We use the IDA characterization of van der Merwe et
al.~\cite{VanderMerwe2021}, as stated in
\cite[Theorem~1]{Berglund2026}: for a memoized state set $M$, IDA
occurs if and only if there are distinct states $p,q$ and a word
$v\in\Sigma^{+}$ with paths
\[
 p\xrightarrow{v}p,
 \qquad p\xrightarrow{v}q,
 \qquad q\xrightarrow{v}q,
\]
such that the displayed cycle at $q$ contains no state of $M$.  Only
that cycle is required to avoid $M$.

As in~\cite{Berglund2026}, positive closure is primitive.  We apply the
Glushkov construction to the displayed syntax without algebraic
simplification, so repeated alternatives are distinct positions.  Let
$x_0$ be the position of the plain $a$, let $x_1,\ldots,x_k$ be the
positions of the $k$ occurrences of $a^{+}$, let $y_b,y_c$ be the
positions of $b,c$, and put
\[
 X_k=\{x_0,x_1,\ldots,x_k\}.
\]
With initial state $q_0$, the nonempty transition values of $G_k$ are
\begin{align*}
 \delta(q_0,a)=\delta(y_b,a)=\delta(y_c,a)&=X_k,\\
 \delta(x_i,b)&=\{y_b\} &&(0\leq i\leq k),\\
 \delta(x_i,c)&=\{y_c\} &&(0\leq i\leq k),\\
 \delta(x_i,a)&=\{x_i\} &&(1\leq i\leq k).
\end{align*}
The final states are $q_0,y_b,y_c$.

Every directed feedback vertex set contains $x_1,\ldots,x_k$, because
each has a self-loop.  After deleting these states, the two cycles
\[
 x_0\longrightarrow y_b\longrightarrow x_0,
 \qquad
 x_0\longrightarrow y_c\longrightarrow x_0
\]
remain.  A single further vertex meets both only if it is $x_0$.
Conversely, deleting $X_k$ leaves only $q_0,y_b,y_c$, with no
transitions among them.  Thus $X_k$ is the unique minimum directed
feedback vertex set.  The Glushkov MFN algorithm returns a minimum
feedback vertex set~\cite[Theorem~4]{Berglund2026}, so
\[
 \MFN(G_k)=X_k.
\]

Memoizing $M_0=\{y_b,y_c\}$ eliminates IDA.  Every nonempty cycle
avoiding $M_0$ is an $a$-self-loop at some $x_i$ with $i\geq1$,
traversed one or more times, and therefore has label $a^r$ for some
$r\geq1$.  If $p\xrightarrow{a^r}p$, then $p=x_j$ for some $j\geq1$,
and the only path on $a^r$ from $x_j$ ends at $x_j$.  It cannot also
end at a distinct state $q$.  Hence $\mu(G_k)\leq2$.

Conversely, let $M$ eliminate IDA and suppose $|M|\leq2$.  If
$y_b\notin M$, then $|X_k|=k+1\geq3$ permits a choice of
$q\in X_k\setminus M$ and $p\in X_k\setminus\{q\}$.  There are paths
\[
 p\xrightarrow{ba}p,
 \qquad p\xrightarrow{ba}q,
 \qquad q\xrightarrow{ba}q.
\]
The last cycle visits only $q$ and $y_b$, neither in $M$, and is an IDA
witness.  Therefore $y_b\in M$.  The same argument with $ca$ gives
$y_c\in M$.  Thus $M_0$ is the unique IDA-eliminating set of size at
most two, and $\mu(G_k)=2$.

Finally, let $M\subsetneq X_k$.  Choose $q\in X_k\setminus M$ and
$p\in X_k\setminus\{q\}$.  The same three paths on $ba$ form an IDA
witness, because the cycle at $q$ visits only $q$ and
$y_b\notin X_k$.  Hence no proper subset of $X_k$ eliminates IDA.
Since deleting $X_k$ is acyclic, $X_k$ does eliminate IDA, and
\[
 \mu_{\mathrm{MFN}}(G_k)=k+1.
\]
Removing any state from the initial MFN set therefore reintroduces
IDA, so $\mathrm{IAR}^{+}$ removes no state, regardless of the ordering.
\end{proof}

The strict inequality
$\mu_{\mathrm{MFN}}(G_k)>\mu(G_k)$ refutes the Glushkov part of
Conjecture~1 in~\cite{Berglund2026}, while the order-independent
$\mathrm{IAR}^{+}$ conclusion refutes Conjecture~2 for Glushkov
automata.  Moreover,
\[
 \mu_{\mathrm{MFN}}(G_k)-\mu(G_k)=k-1,
 \qquad
 \frac{\mu_{\mathrm{MFN}}(G_k)}{\mu(G_k)}=\frac{k+1}{2},
\]
so the additive and multiplicative gaps are unbounded.  The case
$k=2$ gives the explicit five-position counterexample
\[
 \left((a\mid a^{+}\mid a^{+})(b\mid c)\right)^{*}.
\]

\begin{thebibliography}{9}

\bibitem{Berglund2026}
M.~Berglund, B.~van der Merwe, and I.~le Roux,
\newblock Selective memoization for efficient backtracking regular expression matching,
\newblock \emph{Electronic Proceedings in Theoretical Computer Science}
\textbf{446} (2026), 1--17.
\newblock \url{https://doi.org/10.4204/EPTCS.446.1}.

\bibitem{VanderMerwe2021}
B.~van der Merwe, J.~Mouton, S.~van Litsenborgh, and M.~Berglund,
\newblock Memoized regular expressions,
\newblock in \emph{Implementation and Application of Automata (CIAA 2021)},
Lecture Notes in Computer Science, vol.~12803, Springer, 2021,
pp.~39--52.
\newblock \url{https://doi.org/10.1007/978-3-030-79121-6_4}.

\end{thebibliography}

\end{document}